\notePage{Sample: Data Driven Decision Making}{../Materials/2026/3DM_01_Overview_SS26.pdf}{%
\textbf{このページについて}\par
このテンプレートは、ChatGPTがMarkdownからLaTeXへ変換した後の形式例です。実際のノートでは、上部に各スライド画像、下部に日本語訳・解説・試験対策ポイントを入れます。

\medskip
\textbf{数式例}\par
Markdown内の数式は、LaTeXでも次のように保持します。
\[
V_t(S_t)=\max_{x_t\in X_t(S_t)}\mathbb{E}\left[C_t(S_t,x_t,W_{t+1})+V_{t+1}(S_{t+1})\right].
\]

\medskip
\textbf{過去問との関連の書き方}\par
例: SS24 Aufgabe 3cではBellman方程式と、状態・遷移が増えたときに計算不能になる部分が問われています。したがって、この式ではmax、期待値、遷移列挙の意味を説明できる必要があります。
}
